
---

## 4. LaTeX report: 03.mpi.file.transfer.tex

`practicals/03-mpi-file-transfer/03.mpi.file.transfer.tex`

```tex
\documentclass[a4paper,11pt]{article}

\usepackage[margin=2.5cm]{geometry}
\usepackage{listings}
\usepackage{hyperref}

\title{Practical Work 3: MPI File Transfer}
\author{Nguyen Trong Bach \\ Student ID: 23BI14057}
\date{\today}

\begin{document}

\maketitle

\section{Introduction}
This practical extends the TCP file transfer system by using the
Message Passing Interface (MPI). Instead of sockets, processes
communicate via MPI point-to-point operations. The goal is to
implement a simple file transfer system using two MPI processes.

\section{Choice of MPI Implementation}
I chose to use a standard MPI implementation such as MPICH or
OpenMPI because they are widely available, portable and compatible
with the \texttt{mpicc} and \texttt{mpirun} tools provided in the
course environment.

\section{Design of the MPI Service}
The system uses exactly two processes:
\begin{itemize}
  \item Rank 0: acts as the client. It reads the input file and sends
  data chunks to rank 1.
  \item Rank 1: acts as the server. It receives the chunks and writes
  them to the output file.
\end{itemize}

Each chunk transfer consists of two messages: one for the chunk length
and one for the data itself. A chunk length of zero is used as an
End-of-File marker.

\section{System Organization}
The practical is organized in a single C program:
\begin{itemize}
  \item \texttt{mpi\_file\_transfer.c}: contains both client and server
  logic, selected by the MPI rank.
\end{itemize}

The program is compiled with \texttt{mpicc} and executed using
\texttt{mpirun -np 2}.

\section{Implementation Snippet}
Example code from the sender (rank 0):

\begin{lstlisting}[language=C]
while (1) {
    int n = (int)fread(buffer, 1, CHUNK_SIZE, in);
    MPI_Send(&n, 1, MPI_INT, 1, TAG_CHUNK_LEN, MPI_COMM_WORLD);
    if (n == 0) break;
    MPI_Send(buffer, n, MPI_UNSIGNED_CHAR, 1, TAG_CHUNK,
             MPI_COMM_WORLD);
}
\end{lstlisting}

\section{Group Work}
This practical was completed individually.

\section{Conclusion}
By using MPI, the file transfer system no longer depends on TCP
sockets. The communication is handled by MPI sends and receives,
which simplifies some aspects of the networking code and makes it
easier to scale to more processes in the future.

\end{document}

